% --- Archon physics formalization source begin ---
% archon:physics
% archon:covers PhyXMiniProblems/problem_phyx_mini_0137.lean
% archon:source-report reports/phyx_mini/problem_phyx_mini_0137.source.json

\chapter{Physics problem phyx\_mini\_0137}
\label{ch:PhyXMiniProblems_problem_phyx_mini_0137}

\paragraph{Problem source.}
\#\# Physical scenario

A swimmer has dropped her goggles to the bottom of a pool at the shallow end, marked as \textbackslash{}(1.0\textbackslash{},\textbackslash{}mathrm\{m\}\textbackslash{}) deep. But the goggles don't look that deep.And you look straight down into the water.

\#\# Question

How deep do the goggles appear to be?

\#\# Answer choices

- A: A: 0.79m

- B: B: 0.77m

- C: C: 0.75m

- D: D: 0.73m

\#\# Figure caption (auxiliary; use the image as primary evidence)

The image depicts a diagram related to optics, likely involving refraction. Here's a detailed breakdown:

1. **Medium**: The diagram is divided into two parts by a blue horizontal line, representing the interface between water (bottom) and air (top).

2. **Goggles**: At the bottom part, there is an icon labeled "Goggles," which represents an observer's viewpoint under the water.

3. **Vertical Lines**: There are two vertical dashed lines labeled \textbackslash{}( x \textbackslash{}) and \textbackslash{}( x \textbackslash{}) on both sides of the center vertical.

4. **Angles**: 
   - \textbackslash{}( \textbackslash{}theta\_1 \textbackslash{}) is shown at the point of observation with the goggles, representing the angle of the observer's line of sight or light exiting the water.
   - \textbackslash{}( \textbackslash{}theta\_2 \textbackslash{}) and \textbackslash{}( \textbackslash{}theta'\_2 \textbackslash{}) are angles formed by the light paths in the air above the water.

5. **Light Rays**: Two pink lines indicate light rays:
   - One solid pink line extends from the goggles through the water-air interface and into the air.
   - One dashed pink line appears to continue this path above the water, suggesting reflection or refraction, with the two rays diverging as they exit the water.

6. **Distances**:
   - \textbackslash{}( d = 1.0 \textbackslash{}, \textbackslash{}text\{m\} \textbackslash{}) is the vertical distance from the goggles to the water-air interface.
   - \textbackslash{}( d' \textbackslash{}) represents some other vertical distance in the air section, though it is not specifically labeled with a value.

7. **Labels**: The text "Goggles" is beside the goggles icon, indicating the perspective or equipment used for viewing.

Overall, the image likely illustrates the concept of light behavior (refraction and angle of incidence) as it passes from water to air, as seen by someone wearing goggles submerged in water.

\#\# Dataset metadata

Geometrical Optics; Physical Model Grounding Reasoning, Spatial Relation Reasoning

\paragraph{Recorded answer/context.}
C: C: 0.75m

\paragraph{Figure/image path.}
/root/proposal\_for\_physic/science-mango/phyx\_mini\_run/phyx\_data/test\_image/137.png

\paragraph{Formalization target.}
create a compiling Lean file with sorry bodies at `PhyXMiniProblems/problem\_phyx\_mini\_0137.lean`. The Lean declarations must preserve the physical quantities, dimensions or dimensional roles, figure labels, governing-law hypotheses, and final relation expressed by this problem.
Use Mathlib/Physlib names found through LeanExplore where available. If a physics API is missing, introduce faithful local abstractions rather than scalar placeholder aliases.

\begin{theorem}[Physics formalization target]
\label{thm:physics:phyx_mini_0137:target}
\lean{PhyXMiniProblems.ProblemPhyXMini0137.problem_phyx_mini_0137}
\uses{def:physics:phyx-mini-0137:phyxminiproblems-problemphyxmini0137-lengthinmetres, def:physics:phyx-mini-0137:phyxminiproblems-problemphyxmini0137-poolgogglesdiagram, def:physics:phyx-mini-0137:phyxminiproblems-problemphyxmini0137-matchespoolgogglesfigure, def:physics:phyx-mini-0137:phyxminiproblems-problemphyxmini0137-usesstandardwaterairindices, def:physics:phyx-mini-0137:phyxminiproblems-problemphyxmini0137-hasphysicalapparentdepthparameters, def:physics:phyx-mini-0137:phyxminiproblems-problemphyxmini0137-satisfiesparaxialapparentdepthlaws, def:physics:phyx-mini-0137:phyxminiproblems-problemphyxmini0137-isdisplayedapparentdepthanswer}
The assigned autoformalize agent should translate this physics problem into Lean declarations in the covered file, with theorem and lemma proof bodies written as `by sorry`.
\end{theorem}
\begin{proof}
This is an autoformalization task, not a proof task. Produce statements that can later be proved by the physics prover without weakening the physical model.
\end{proof}

% --- Archon declaration topology begin ---
\section{Declaration topology}
\begin{definition}[Dim Length]
  \label{def:physics:phyx-mini-0137:phyxminiproblems-problemphyxmini0137-dimlength}
  \lean{PhyXMiniProblems.ProblemPhyXMini0137.DimLength}
  A physical length, represented independently of any choice of unit
\end{definition}
\begin{proof}
  This is the declared model definition of Dim Length.
\end{proof}
\begin{definition}[length Readout]
  \label{def:physics:phyx-mini-0137:phyxminiproblems-problemphyxmini0137-lengthreadout}
  \lean{PhyXMiniProblems.ProblemPhyXMini0137.lengthReadout}
  \uses{def:physics:phyx-mini-0137:phyxminiproblems-problemphyxmini0137-dimlength}
  The scalar readout of a physical length in the selected length unit
\end{definition}
\begin{proof}
  This is the declared model definition of length Readout.
\end{proof}
\begin{definition}[length In Metres]
  \label{def:physics:phyx-mini-0137:phyxminiproblems-problemphyxmini0137-lengthinmetres}
  \lean{PhyXMiniProblems.ProblemPhyXMini0137.lengthInMetres}
  \uses{def:physics:phyx-mini-0137:phyxminiproblems-problemphyxmini0137-dimlength, def:physics:phyx-mini-0137:phyxminiproblems-problemphyxmini0137-lengthreadout}
  The scalar readout of a physical length in metres
\end{definition}
\begin{proof}
  This is the declared model definition of length In Metres.
\end{proof}
\begin{definition}[Optical Medium]
  \label{def:physics:phyx-mini-0137:phyxminiproblems-problemphyxmini0137-opticalmedium}
  \lean{PhyXMiniProblems.ProblemPhyXMini0137.OpticalMedium}
  The homogeneous optical media on the two sides of the pool surface
\end{definition}
\begin{proof}
  This is the declared model definition of Optical Medium.
\end{proof}
\begin{definition}[Ray Path Label]
  \label{def:physics:phyx-mini-0137:phyxminiproblems-problemphyxmini0137-raypathlabel}
  \lean{PhyXMiniProblems.ProblemPhyXMini0137.RayPathLabel}
  Labels for the physical and apparent ray paths drawn in the figure
\end{definition}
\begin{proof}
  This is the declared model definition of Ray Path Label.
\end{proof}
\begin{definition}[Figure Angle Label]
  \label{def:physics:phyx-mini-0137:phyxminiproblems-problemphyxmini0137-figureanglelabel}
  \lean{PhyXMiniProblems.ProblemPhyXMini0137.FigureAngleLabel}
  The three angle labels printed in the figure
\end{definition}
\begin{proof}
  This is the declared model definition of Figure Angle Label.
\end{proof}
\begin{definition}[Viewing Regime]
  \label{def:physics:phyx-mini-0137:phyxminiproblems-problemphyxmini0137-viewingregime}
  \lean{PhyXMiniProblems.ProblemPhyXMini0137.ViewingRegime}
  Whether the viewing axis is normal or oblique to the water surface
\end{definition}
\begin{proof}
  This is the declared model definition of Viewing Regime.
\end{proof}
\begin{definition}[Pool Goggles Diagram]
  \label{def:physics:phyx-mini-0137:phyxminiproblems-problemphyxmini0137-poolgogglesdiagram}
  \lean{PhyXMiniProblems.ProblemPhyXMini0137.PoolGogglesDiagram}
  \uses{def:physics:phyx-mini-0137:phyxminiproblems-problemphyxmini0137-dimlength, def:physics:phyx-mini-0137:phyxminiproblems-problemphyxmini0137-opticalmedium, def:physics:phyx-mini-0137:phyxminiproblems-problemphyxmini0137-figureanglelabel, def:physics:phyx-mini-0137:phyxminiproblems-problemphyxmini0137-viewingregime}
  The dimensionful distances and dimensionless optical readouts appearing in the water--air apparent-depth diagram. lateralOffset is the common figure label x; realDepth and apparentDepth are the labels d and d'. No numerical value of apparentDepth is stored in this setup
\end{definition}
\begin{proof}
  This is the declared model definition of Pool Goggles Diagram.
\end{proof}
\begin{definition}[Matches Pool Goggles Figure]
  \label{def:physics:phyx-mini-0137:phyxminiproblems-problemphyxmini0137-matchespoolgogglesfigure}
  \lean{PhyXMiniProblems.ProblemPhyXMini0137.MatchesPoolGogglesFigure}
  \uses{def:physics:phyx-mini-0137:phyxminiproblems-problemphyxmini0137-lengthinmetres, def:physics:phyx-mini-0137:phyxminiproblems-problemphyxmini0137-poolgogglesdiagram}
  Readouts supplied by the problem and diagram: d = 1.0 m, normal viewing, and equality of the air-ray angle θ₂ with the corresponding angle θ₂' of its backward extension relative to the parallel vertical normals
\end{definition}
\begin{proof}
  This is the declared model definition of Matches Pool Goggles Figure.
\end{proof}
\begin{definition}[Uses Standard Water Air Indices]
  \label{def:physics:phyx-mini-0137:phyxminiproblems-problemphyxmini0137-usesstandardwaterairindices}
  \lean{PhyXMiniProblems.ProblemPhyXMini0137.UsesStandardWaterAirIndices}
  \uses{def:physics:phyx-mini-0137:phyxminiproblems-problemphyxmini0137-poolgogglesdiagram}
  Standard dimensionless refractive-index data used for water and air
\end{definition}
\begin{proof}
  This is the declared model definition of Uses Standard Water Air Indices.
\end{proof}
\begin{definition}[Has Physical Apparent Depth Parameters]
  \label{def:physics:phyx-mini-0137:phyxminiproblems-problemphyxmini0137-hasphysicalapparentdepthparameters}
  \lean{PhyXMiniProblems.ProblemPhyXMini0137.HasPhysicalApparentDepthParameters}
  \uses{def:physics:phyx-mini-0137:phyxminiproblems-problemphyxmini0137-lengthinmetres, def:physics:phyx-mini-0137:phyxminiproblems-problemphyxmini0137-poolgogglesdiagram}
  Positivity and acute-angle conditions selecting the physical ray branch
\end{definition}
\begin{proof}
  This is the declared model definition of Has Physical Apparent Depth Parameters.
\end{proof}
\begin{definition}[Satisfies Exact Water Air Snell Law]
  \label{def:physics:phyx-mini-0137:phyxminiproblems-problemphyxmini0137-satisfiesexactwaterairsnelllaw}
  \lean{PhyXMiniProblems.ProblemPhyXMini0137.SatisfiesExactWaterAirSnellLaw}
  \uses{def:physics:phyx-mini-0137:phyxminiproblems-problemphyxmini0137-poolgogglesdiagram}
  Exact Snell law for the water-to-air interface, n\_water sin θ₁ = n\_air sin θ₂. The numerical theorem below uses its paraxial first-order form, stated separately, because the displayed multiple-choice answer is the standard paraxial apparent-depth value
\end{definition}
\begin{proof}
  This is the declared model definition of Satisfies Exact Water Air Snell Law.
\end{proof}
\begin{definition}[Satisfies Paraxial Apparent Depth Laws]
  \label{def:physics:phyx-mini-0137:phyxminiproblems-problemphyxmini0137-satisfiesparaxialapparentdepthlaws}
  \lean{PhyXMiniProblems.ProblemPhyXMini0137.SatisfiesParaxialApparentDepthLaws}
  \uses{def:physics:phyx-mini-0137:phyxminiproblems-problemphyxmini0137-lengthreadout, def:physics:phyx-mini-0137:phyxminiproblems-problemphyxmini0137-poolgogglesdiagram}
  First-order straight-down ray optics. The first field is linearized Snell law; the other two fields are the small-angle ray triangles x = d θ₁ and x = d' θ₂'. They are required in every length unit, making their dimensional content explicit and independent of a chosen readout
\end{definition}
\begin{proof}
  This is the declared model definition of Satisfies Paraxial Apparent Depth Laws.
\end{proof}
\begin{lemma}[apparent depth index relation]
  \label{lem:physics:phyx-mini-0137:phyxminiproblems-problemphyxmini0137-apparent-depth-index-relation}
  \lean{PhyXMiniProblems.ProblemPhyXMini0137.apparent_depth_index_relation}
  \uses{def:physics:phyx-mini-0137:phyxminiproblems-problemphyxmini0137-lengthreadout, def:physics:phyx-mini-0137:phyxminiproblems-problemphyxmini0137-poolgogglesdiagram, def:physics:phyx-mini-0137:phyxminiproblems-problemphyxmini0137-matchespoolgogglesfigure, def:physics:phyx-mini-0137:phyxminiproblems-problemphyxmini0137-hasphysicalapparentdepthparameters, def:physics:phyx-mini-0137:phyxminiproblems-problemphyxmini0137-satisfiesparaxialapparentdepthlaws}
  Eliminating the nonzero ray angles and the common offset gives the paraxial apparent-depth relation n\_water d' = n\_air d in any length unit
\end{lemma}
\begin{proof}
  The conclusion follows from the governing relations and figure readouts named in the statement of apparent depth index relation.
\end{proof}
\begin{definition}[Answer Choice]
  \label{def:physics:phyx-mini-0137:phyxminiproblems-problemphyxmini0137-answerchoice}
  \lean{PhyXMiniProblems.ProblemPhyXMini0137.AnswerChoice}
  Labels of the four displayed multiple-choice answers
\end{definition}
\begin{proof}
  This is the declared model definition of Answer Choice.
\end{proof}
\begin{definition}[displayed Apparent Depth Metres]
  \label{def:physics:phyx-mini-0137:phyxminiproblems-problemphyxmini0137-displayedapparentdepthmetres}
  \lean{PhyXMiniProblems.ProblemPhyXMini0137.displayedApparentDepthMetres}
  \uses{def:physics:phyx-mini-0137:phyxminiproblems-problemphyxmini0137-answerchoice}
  Apparent depth in metres printed next to each answer choice
\end{definition}
\begin{proof}
  This is the declared model definition of displayed Apparent Depth Metres.
\end{proof}
\begin{definition}[Is Displayed Apparent Depth Answer]
  \label{def:physics:phyx-mini-0137:phyxminiproblems-problemphyxmini0137-isdisplayedapparentdepthanswer}
  \lean{PhyXMiniProblems.ProblemPhyXMini0137.IsDisplayedApparentDepthAnswer}
  \uses{def:physics:phyx-mini-0137:phyxminiproblems-problemphyxmini0137-lengthinmetres, def:physics:phyx-mini-0137:phyxminiproblems-problemphyxmini0137-poolgogglesdiagram, def:physics:phyx-mini-0137:phyxminiproblems-problemphyxmini0137-answerchoice, def:physics:phyx-mini-0137:phyxminiproblems-problemphyxmini0137-displayedapparentdepthmetres}
  A displayed answer agrees with the apparent-depth readout of the model
\end{definition}
\begin{proof}
  This is the declared model definition of Is Displayed Apparent Depth Answer.
\end{proof}
% --- Archon declaration topology end ---
% --- Archon physics formalization source end ---
